% Lösungen Einheit 1 von 2 – Verteilung aufstellen: die Werte der Zufallsgröße aus den Regeln gewinnen (alle Ergebnisfolgen durchrechnen), die Tabelle durch Abzählen füllen (günstige Paare je Wert), fehlende Wahrscheinlichkeiten über die Summe eins (zusammenbau v0.1)
\einheitenkopf{Einheit 1 von 2 · Verteilung aufstellen: die Werte der Zufallsgröße aus den Regeln gewinnen (alle Ergebnisfolgen durchrechnen), die Tabelle durch Abzählen füllen (günstige Paare je Wert), fehlende Wahrscheinlichkeiten über die Summe eins}

\erg{6}{a) alle vier Einträge der unteren Zeile – der fehlende bei $k = 1$ ergänzt sie zu eins}
\erg{7}{a) $0, 1, 2$ – keinmal, einmal oder zweimal Kopf \quad b) $\frac{1}{16}$; $\frac{2}{16}$; $\frac{3}{16}$; $\frac{4}{16}$ \quad c) $-1$, $1$, $7$: Die Folgen 1-1, 1-3, 3-1, 3-3 liefern die Produkte 1, 3, 3, 9; minus Einsatz ergibt das $-1$, $1$, $1$, $7$. \quad d) $2p + 5p + p = 8p = 1$, also $p = \frac{1}{8}$ – die Beträge spielen dabei keine Rolle \quad e) $P(X = 2) = 0$, denn stecken zwei richtig, steckt auch der dritte richtig; $P(X = 3) = 1 - \frac{1}{3} - \frac{1}{2} = \frac{1}{6}$}
\erg{8}{a) Fehler: 2 und 5 sowie 5 und 2 liefern denselben Wert, er wird einmal gelistet und doppelt gezählt. Richtig: Werte 4, 7, 10 mit $P(X = 7) = \frac{1}{2}$ \quad b) Die Werte teilen alle Ergebnisse vollständig und ohne Überschneidung auf; zusammen bilden sie das sichere Ereignis. \quad c) Gewinn $9$ € mit $\frac{1}{16}$, $0$ € mit $\frac{6}{16}$, $-1$ € mit $\frac{9}{16}$; der Einsatz ist mit $P = \frac{9}{16}$ verloren \quad d) Säulen der Höhen $0{,}125$; $0{,}375$; $0{,}375$; $0{,}125$ über $0$ bis $3$ – symmetrisch zur Mitte.}

8d)

\saeulen{0/0.125,1/0.375,2/0.375,3/0.125}{0.5}{0.125}{$P(X = k)$}

